\exercisegroup

\begin{enumerate}
\def\labelenumi{\arabic{enumi}.}
\item
  For a graph \(\mathbf{G}\) to be the graph of an equivalence relation on a set \(\mathbf{E}\) , it is necessary and sufficient that \(\mathrm{pr}_1\mathbf{G} = \mathbf{E}\) , \(\mathbf{G} \circ \overline{\mathbf{G}} \circ \mathbf{G} = \mathbf{G}\) , and \(\Delta_{\mathbf{E}} \subset \mathbf{G}\) ( \(\Delta_{\mathbf{E}}\) being the diagonal of \(\mathbf{E}\) ).
\item
  If \(\mathbf{G}\) is a graph such that
\end{enumerate}

\[
\mathbf {G} \circ \overline {{{{\mathbf {G}}}}} ^ {- 1} \circ \mathbf {G} = \mathbf {G},
\]

show that \(\overline{\mathbf{G}}\circ \mathbf{G}\) and \(\mathbf{G}\circ \overline{\mathbf{G}}\) are graphs of equivalences on \(\mathfrak{pr}_1\mathbf{G}\) and \(\mathfrak{pr}_2\mathbf{G}\) respectively.

\begin{enumerate}
\def\labelenumi{\arabic{enumi}.}
\setcounter{enumi}{2}
\item
  Let E be a set, A a subset of E, and R the equivalence relation associated with the mapping \(X \rightarrow X \cap A\) of \(\mathfrak{P}(E)\) into \(\mathfrak{P}(E)\) . Show that there exists a bijection of \(\mathfrak{P}(A)\) onto the quotient set \(\mathfrak{P}(E)/R\) .
\item
  Let \(G\) be the graph of an equivalence on a set \(E\) . Show that if \(A\) is a graph such that \(A \subset G\) and \(\operatorname{pr}_1 A = E\) (resp. \(\operatorname{pr}_2 A = E\) ), then
\end{enumerate}

\[
\mathbf {G} \circ \mathbf {A} = \mathbf {G} (\text { resp. } \quad \mathbf {A} \circ \mathbf {G} = \mathbf {G});
\]

furthermore, if B is any graph, we have \((G \cap B) \circ A = G \cap (B \circ A)\) (resp. \(A \circ (G \cap B) = G \cap (A \circ B)\) ).

\begin{enumerate}
\def\labelenumi{\arabic{enumi}.}
\setcounter{enumi}{4}
\item
  Show that every intersection of graphs of equivalences on a set E is the graph of an equivalence on E. Give an example of two equivalences on a set E such that the union of their graphs is not the graph of an equivalence on E.
\item
  Let G, H be the graphs of two equivalences on E. Then \(G \circ H\) is the graph of an equivalence on E if and only if \(G \circ H = H \circ G\) . The graph \(G \circ H\) is then the intersection of all the graphs of equivalences on E which contain both G and H.
\item
  Let \(G_{0}\) , \(G_{1}\) , \(H_{0}\) , \(H_{1}\) be the graphs of four equivalences on a set E such that \(G_{1} \cap H_{0} = G_{0} \cap H_{1}\) and \(G_{1} \circ H_{0} = G_{0} \circ H_{1}\) . For each \(x \in E\) , let \(R_{0}\) (resp. \(S_{0}\) ) be the relation induced on \(\mathbf{G}_{1}(x)\) (resp. \(H_{1}(x)\) ) by the equivalence relation \((x, y) \in \mathbf{G}_{0}\) (resp. \((x, y) \in \mathbf{H}_{0}\) ). Show that there exists a bijection of the quotient set \(\mathbf{G}_{1}(x)/\mathbf{R}_{0}\) onto the quotient set
\end{enumerate}

\[
\mathrm{H} _ {1} (x) / \mathrm{S} _ {0}.
\]

(If \(\mathbf{A} = \mathbf{G}_1(x) \cap \mathbf{H}_1(x)\) , show that both quotient sets are in one-to-one correspondence with the quotient set of \(\mathbf{A}\) by the equivalence relation induced by \(\mathbf{R}_0\) on \(\mathbf{A}\) ; this relation is equivalent to that induced by \(\mathbf{S}_0\) on \(\mathbf{A}\) .)

\begin{enumerate}
\def\labelenumi{\arabic{enumi}.}
\setcounter{enumi}{7}
\item
  Let E, F be two sets, let R be an equivalence relation on F, and let f be a mapping of E into F. If S is the equivalence relation which is the inverse image of R under f, and if \(A = f \langle E \rangle\) , define a canonical bijection of E/S onto \(A/R_{A}\) .
\item
  Let F, G be two sets, let R be an equivalence relation on F, let p be the canonical mapping of F onto F/R, and let f be a surjection of G onto F/R. Show that there exists a set E, a surjection g of E onto F, and a surjection h of E onto G, such that \(p \circ g = f \circ h\) .
\item
  \begin{enumerate}
  \def\labelenumii{(\alph{enumii})}
  \tightlist
  \item
    If \(R\{x, y\}\) is any relation, then `` \(R\{x, y\}\) and \(R\{y, x\}\) '' is a symmetric relation. Under what condition is it reflexive on a set E?
  \end{enumerate}
\end{enumerate}

*(b) Let \(\mathbb{R}\{x, y\}\) be a reflexive and symmetric relation on a set \(\mathbf{E}\) . Let \(\mathbb{S}\{x, y\}\) be the relation ``there exists an integer \(n > 0\) and a sequence \((x_i)_{0 \leqslant i \leqslant n}\) of elements of \(\mathbf{E}\) such that \(x_0 = x\) , \(x_n = y\) , and for each index \(i\) such that \(0 \leqslant i < n\) , \(\mathbb{R}\{x_i, x_{i+1}\}\) ''. Show that \(\mathbb{S}\{x, y\}\) is an equivalence relation on \(\mathbf{E}\) and that its graph is the smallest of all graphs of equivalences on \(\mathbf{E}\) which contain the graph of \(\mathbb{R}\) . The equivalence classes with respect to \(\mathbb{S}\) are called the connected components of \(\mathbf{E}\) with respect to the relation \(\mathbb{R}\) .

\begin{enumerate}
\def\labelenumi{(\alph{enumi})}
\setcounter{enumi}{2}
\tightlist
\item
  Let \(\mathfrak{F}\) be the set of subsets A of E such that, for each pair of elements \((y,z)\) such that \(y\in A\) and \(z\in E - A\) , we have ``not R \(\{y,z\} "\) .
\end{enumerate}

For each \(x \in \mathbf{E}\) , show that the intersection of the sets \(\mathbf{A} \in \mathfrak{F}\) such that \(x \in \mathbf{A}\) is the connected component of \(x\) with respect to the relation \(\mathbf{R}_{*}\) .

\begin{enumerate}
\def\labelenumi{\arabic{enumi}.}
\setcounter{enumi}{10}
\tightlist
\item
  \begin{enumerate}
  \def\labelenumii{(\alph{enumii})}
  \tightlist
  \item
    Let \(\mathbb{R}\{x, y\}\) be a reflexive and symmetric relation on a set \(\mathbf{E}\) . \(\mathbf{R}\) is said to be intransitive of order 1 if for any four distinct elements \(x, y, z, t\) of \(\mathbf{E}\) the relations \(\mathbb{R}\{x, y\}, \mathbb{R}\{x, z\}, \mathbb{R}\{x, t\}, \mathbb{R}\{y, z\}\) , and \(\mathbb{R}\{y, t\}\) imply \(\mathbb{R}\{z, t\}\) . A subset \(\mathbf{A}\) of \(\mathbf{E}\) is said to be stable with respect to the relation \(\mathbb{R}\) if \(\mathbb{R}\{x, y\}\) for all \(x\) and \(y\) in \(\mathbf{A}\) . If \(a\) and \(b\) are two distinct elements of \(\mathbf{E}\) such that \(\mathbb{R}\{a, b\}\) , show that the set \(\mathbf{C}(a, b)\) of elements \(x \in \mathbf{E}\) such that \(\mathbb{R}\{a, x\}\) and \(\mathbb{R}\{b, x\}\) is stable and that
  \end{enumerate}
\end{enumerate}

\[
\mathbf {C} (x, y) = \mathbf {C} (a, b)
\]

for each pair of distinct elements \(x, y\) of \(C(a, b)\) . The sets \(C(a, b)\) (for each ordered pair \((a, b)\) such that \(R\{a, b\}\) ) and the connected components (Exercise 10) with respect to \(R\) which consist of a single element are called the constituents of \(E\) with respect to the relation \(R\) . Show that the intersection of two distinct constituents of \(E\) contains at most one element and that if \(A, B, C\) are three mutually distinct constituents, at least one of the sets \(A \cap B, B \cap C, C \cap A\) is empty.

\begin{enumerate}
\def\labelenumi{(\alph{enumi})}
\setcounter{enumi}{1}
\item
  Conversely, let \((\mathbf{X}_{\lambda})_{\lambda \in \mathbf{L}}\) be a covering of a set \(\mathbf{E}\) consisting of non-empty subsets of \(\mathbf{E}\) and having the following properties: (1) if \(\lambda, \mu\) are two distinct indices, \(\mathbf{X}_{\lambda} \cap \mathbf{X}_{\mu}\) contains at most one element; (2) if \(\lambda, \mu, \nu\) are three distinct indices, then at least one of the three sets \(\mathbf{X}_{\lambda} \cap \mathbf{X}_{\mu}, \mathbf{X}_{\mu} \cap \mathbf{X}_{\nu}, \mathbf{X}_{\nu} \cap \mathbf{X}_{\lambda}\) is empty. Let \(\mathbf{R}\{x, y\}\) be the relation ``there exists \(\lambda \in \mathbf{L}\) such that \(x \in \mathbf{X}_{\lambda}\) and \(y \in \mathbf{X}_{\lambda}\) ; show that \(\mathbf{R}\) is reflexive on \(\mathbf{E}\) , symmetric and intransitive of order 1, and that the \(\mathbf{X}_{\lambda}\) are the constituents of \(\mathbf{E}\) with respect to \(\mathbf{R}\) .
\item
  * Similarly, a relation \(R\{x, y\}\) which is reflexive and symmetric on E is said to be intransitive of order n-3 if, for every family \((x_{i})_{1 \leqslant i \leqslant n}\) of distinct elements of E, the relations \(R\{x_{i}, x_{j}\}\) for each pair \((i, j) \neq (n-1, n)\) imply \(R\{x_{n-1}, x_{n}\}\) . Generalize the results of (a) and (b) to intransitive relations of arbitrary order. Show that a relation which is intransitive of order p is also intransitive of order q for all q \textgreater{} p.~*
\end{enumerate}
